\section{Introduction}
\label{sec:intro}

An inclusion proof lets a verifier, given only a short commitment to a large
dataset, confirm that a particular element is really part of that dataset --
without holding the dataset itself. The standard tool for this is the Merkle
tree \citep{merkle1980}: leaves are hashed, pairs of hashes are hashed
together recursively up to a root, and an inclusion proof consists of the
$O(\log N)$ sibling hashes on the path from a leaf to the root. Both proof
size and verification time are $O(\log N)$ in the number of elements $N$.

This efficiency comes at a structural cost: the committed data must be hashed
\emph{as a tree}. Many real systems, however, do not have that freedom. A
Bitcoin transaction's witness data, for instance, can carry an optional
\emph{annex}, whose commitment is defined by BIP-341 \citep{bip341} as a
single, ordinary SHA-256 of the annex bytes -- a flat, linear
Merkle--Damg\aa rd hash, with no tree structure available at the consensus
layer: consensus only ever computes $\mathrm{SHA256}(\text{annex bytes})$,
once, linearly.

\subsection*{Threat model}

The setting this paper targets is adversarial on the committing side. One
party -- the \emph{committer} -- publishes data and its digest; other
parties -- \emph{opponents} -- later need to convince a verifier who holds
only the digest (a Bitcoin script, a dispute program, a contract) of facts
about that data. The requirement is that, for \emph{every} byte string the
committer might publish, an honest opponent who holds it can always prove
something in $O(\log N)$ space and verification time: either the
inclusion of any element, or -- if the commitment is malformed so that
inclusion proofs for the published data cannot be built -- the
malformation itself. A malicious committer will choose whichever of the
two is weaker, so if either proof costs $O(N)$, the scheme as a whole
offers only $O(N)$ guarantees.

This is what rules out the obvious way of embedding a Merkle tree in a
linearly hashed field. Appending a Merkle root as the field's last 32 bytes
gives logarithmic inclusion proofs (one compression call from a revealed
midstate opens the root, then an ordinary Merkle path), but a committer can
publish a root that disagrees with the published data at a single position
near the start. Exposing that requires proving which bytes physically sit at
that position, and in a linear hash this means recomputing the chain from
there to the end: $O(N)$. Section~\ref{sec:annex-alternatives} compares this
and other alternatives in detail.

This raises a natural question: can a \emph{linear} hash stream support
$O(\log N)$ inclusion proofs \emph{and} $O(\log N)$ proofs of malformation,
without changing what is actually hashed at the commitment layer?

\subsection*{This paper}

We show that it can. The idea is simple to state: periodically, as the
stream is hashed, re-inject the hash function's own current internal state
back into the stream as if it were ordinary message content. Later, a proof
can reveal one of these re-injected states instead of the raw data it
summarizes, letting the verifier skip an entire sub-range of the file in a
single hash-function call. Applied recursively at doubling distances
($2, 4, 8, 16, \dots$), this yields the same $O(\log N)$ proof size and
$O(\log N)$ verification time as a Merkle tree -- while the file, from the
outside, is still just one ordinary flat hash of one ordinary byte string.
Because the re-injected states make any position of the stream reachable in
$O(\log N)$ calls, the same mechanism also lets an opponent prove, in
$O(\log N)$, that a published file deviates from the canonical layout
(Section~\ref{sec:fraud-proofs}), meeting both halves of the threat model.
The price is size: the structured file is about twice the size of the raw
data (Section~\ref{sec:construction}).
We call the resulting structure \textbf{FlatMerkle}.

The security of the construction rests on a specific, well-studied but
non-standard assumption: \emph{free-start} collision resistance of the
hash function's raw compression step (Section~\ref{sec:security}) -- the
hardness of finding a compression-function collision when the adversary is
free to choose the starting state, not just the message. This is strictly
stronger than the ordinary (IV-fixed) collision resistance that hash
functions are normally analyzed under, and we state it as an explicit,
separate assumption rather than treating it as automatic.

\subsection*{Contributions}

\begin{itemize}
\item A formal recursive construction of the structured file, and a
  closed-form formula, $\mathrm{position}(a) = 2a - \popcount(a)$, that
  locates any original element's \emph{item index} within the structured
  stream without a lookup table -- a count of items, not literally a byte
  offset; multiplying by the fixed per-item byte width gives the byte
  offset (Section~\ref{sec:construction}). The inclusion-proof path itself
  comes from a separate, simpler branch rule (Section~\ref{sec:algorithm}).
\item Single-element proof construction and verification algorithms, with
  worked examples and illustrations for a 16-element file
  (Sections~\ref{sec:algorithm}--\ref{sec:examples}).
\item A security reduction: forging a FlatMerkle inclusion proof, or making
  a single digest accept two conflicting answers for the same index,
  implies finding a free-start collision in the underlying compression
  function. This binding guarantee does \emph{not} extend to files that
  were never canonically built in the first place -- we give a working
  forgery against that stronger claim -- so we also give a complementary
  $O(N)$-detect/$O(\log N)$-prove fraud-proof mechanism that lets any party
  holding the full file catch and prove a non-canonical layout, turning
  soundness against a malicious publisher into an optimistic rather than
  unconditional guarantee (Section~\ref{sec:security}).
\item Optimal element-size analysis minimizing average proof size, both for
  a size-flexible toy primitive and for real, fixed-block SHA-256
  (Section~\ref{sec:optimal-size}).
\item A batched multi-element inclusion proof construction with measured
  space savings against $m$ independent proofs (Section~\ref{sec:multiproof}).
\item A measured construction-speed comparison against Merkle trees, using
  an identical from-scratch SHA-256 compression function on both sides
  (Section~\ref{sec:benchmark}).
\item A concrete application: committing a FlatMerkle structured file inside
  a Bitcoin transaction annex, so that BIP-341's existing
  $\mathrm{SHA256}(\text{annex})$ commitment doubles as a FlatMerkle root,
  with no consensus changes (Section~\ref{sec:bitcoin-annex}).
\item A design sketch composing this annex application with BitVMX and
  BitVM3, two Bitcoin computation-verification protocols, and ESSPI
  \citep{lerner2024bitvmx,linuswoll2026bitvm3,lerner2025esspi} to use a
  FlatMerkle-committed annex as the data-availability layer for an
  optimistic rollup, with watchtowers disputing both non-canonical layouts
  and individually incorrect state transitions
  (Section~\ref{sec:rollup-da}).
\item A from-scratch gas-cost comparison of four ways to commit rollup
  calldata on-chain, using current protocol constants, on Ethereum,
  Rootstock, and Arbitrum
  \citep{eip2028,eip4844,eip7623,eip2929,eip2200,eip152,rsk2022,rskip400,rskip153,rskip239,rskjgascost,arbos40,arbos51}:
  the three chains disagree on whether an in-contract FlatMerkle
  commitment or an in-contract Keccak Merkle tree is cheaper, tracing the
  difference to two separable protocol design choices (a calldata floor
  and a warm-call discount for precompiles) rather than to anything about
  FlatMerkle itself; both remain a clear improvement over submitting an
  externally built FlatMerkle file or persisting the raw data in storage,
  on all three chains (Section~\ref{sec:rollup-eth}).
\item A comparison against the closest prior art we identified, authenticated
  append-only skip lists \citep{maniatis2003,goodrich2001}
  (Section~\ref{sec:related-work}).
\end{itemize}

Every number, table, and figure about the construction itself was
produced by an executable, self-tested reference implementation (a
size-flexible toy compression function and a real-SHA-256 variant), not
asserted from hand calculation. The exceptions are the EVM gas-cost figures
of Section~\ref{sec:rollup-eth} and Appendix~\ref{app:eth-da}, which are
derived by hand from cited protocol constants plus a flagged estimate, and
the build-time extrapolation of Section~\ref{sec:benchmark}. For single-element proofs, the self-tests flip every
individual revealed value, one at a time, for every index at every tested
$N$, and confirm rejection; multiproofs are checked the same way across
many randomized target sets, tampering the data, the claimed data, one
shared reveal, the index-to-pair assignment, and the pair set per trial. The reference implementation -- filenames retain ``LogLink,'' this
project's working name before it was renamed FlatMerkle; the construction
itself is unaffected -- \texttt{loglink.py} (toy compression function),
\texttt{loglink\_annex.py} (real SHA-256),
\texttt{illustrate.py} (figures), \texttt{fraud\_proof.py} and
\texttt{simulate\_fraud\_proofs.py} (Section~\ref{sec:fraud-proofs}),
\texttt{experiments.py} and \texttt{benchmark.py} (Sections~\ref{sec:multiproof}
and~\ref{sec:benchmark}) -- accompanies this paper; a line-by-line Java
transliteration of the core algorithms also exists for cross-language
sanity-checking, but should not be read as an independently designed
second implementation.
